\chapter{Validation on the USD LIBOR-to-SOFR Transition}\label{ch6-usd-validation}\label{ch:usd}

\begin{quote}\itshape
This chapter specifies the empirical test of the thesis's two results on the one transition whose answer is already known. The identification result of Chapter~\ref{ch:identifiability} is applied before the successor rate had an options market and checked once its smile became observable. The certificate of Chapter~\ref{ch:certification} is deployed at the market convention, with realised coverage reported against each deficit term separately. The blocked and unblocked routes of Chapter~\ref{ch:sharpness} are compared on equal footing at the record lengths available. Every criterion precedes the data, and the grid cells on which a criterion cannot be informative are counted in advance. No number here is computed from market data: the block-layout, deficit and sizing arithmetic uses only the record length, the level $\alpha$ and the assumed mixing class.
\end{quote}

\section{The transition of 2021 to 2023 as a natural experiment}\label{sec:usd-experiment}

\subsection{The design in one paragraph}

Chapter~\ref{ch:identifiability} makes a prediction that can be checked only where a predecessor benchmark has been replaced by a successor and where both have had an options market, one after the other. The transition from three-month \gls{libor} in United States dollars to compounded \gls{sofr} plus a fixed spread is such a case, and the only one available at the scale the thesis needs. The spread was fixed on a single date, 5~March 2021, and was constant from then on; the predecessor option market was deep and continuously quoted until the end of 2021 and was extinguished with the panel on 30~June 2023; the successor option market did not exist at the start of that interval and was quoted by the end of it. On any date in between an analyst could form the prediction from the predecessor smile, the two curves and the published spread, and could not have formed it from successor quotes because there were none. The comparison of that prediction with the observation that followed is the experiment.

What makes it an experiment rather than an in-sample exercise is that the concealment was not the analyst's choice: in a backtest one withholds data that exist, here they did not exist. The date on which the successor surface began to be quoted was set by dealer behaviour, clearing-house rules and a regulatory timetable, none responsive to any prediction of this kind. That exogeneity is the whole of the design, in the weak sense of an exogenous timing of observability and no stronger.

\subsection{What the experiment can settle}

Four questions are answerable on this record, and each is already a criterion of Section~\ref{sec:ident-validation} or of Section~\ref{sec:usd-cert} below. Whether the identified interval of Theorem~\ref{thm:ident-two} contains the successor at-the-money volatility the market later quoted, which is criterion~V1, with criterion~V5 separating a wrong model class from an underestimated basis volatility. Whether the interval is narrow enough to be of use, which is criterion~V3, a threshold that Chapter~\ref{ch:identifiability}'s own illustration fails. Whether the market convention's answer is wrong in the direction and by the magnitude the theory predicts, which is criterion~V4 and the only criterion comparing the thesis's construction with what a desk would otherwise do. And whether the certificate of Chapter~\ref{ch:certification} delivers out of sample the coverage it states, with the policy fixed on one part of the record, the threshold computed on a second and the breaches counted on a third.

\subsection{What the experiment cannot settle}

Five things are outside its reach, each a limit of the design and not of the execution, and each treated at length in Section~\ref{sec:usd-threats} rather than here.

It cannot settle the coverage probability of the identified interval, because the interval has none (Section~\ref{sec:ident-validation}, ``What the test cannot show''); the containment rates of Section~\ref{sec:usd-ident} are descriptive frequencies over one realised path. It cannot settle whether the basis carries a volatility risk premium, the nuisance parameters being objects under the pricing measure and every available estimate an object under the physical measure; criterion~V5 measures the agreement of the two and no lengthening of any window reduces the gap. It cannot, on containment evidence alone, distinguish an interval that is correct from one that is too wide, which is what criterion~V3 is for. It cannot settle anything about the South African market by transport, nor rule out circularity where the vendor exposes no quoted-or-derived flag; the first is why Chapter~\ref{ch:zar} takes no parameter estimate from here and the second is why the windows fall where they do.

And it cannot, with one transition, establish that any of this generalises. The unit of replication is the transition, and there is one; within it the unit is the observation-date-and-expiry pair, of which there are many but which are strongly dependent, so every statistic below carries a dependence correction stated beside it.

\subsection{Map of the chapter}

Section~\ref{sec:usd-protocol} fixes the protocol, Section~\ref{sec:usd-ident} the identification test, Section~\ref{sec:usd-cert} the certificate test, Section~\ref{sec:usd-methods} the comparison methods, Section~\ref{sec:usd-routes} the route comparison, Section~\ref{sec:usd-falsify} the falsification rules, Section~\ref{sec:usd-threats} the threats to validity and Section~\ref{sec:usd-discussion} the discussion.

Every dataset named below is described in Appendix~\ref{app:data} and nowhere else; this chapter gives only the analysis and, with each dataset, a note saying what must be pulled and from where.

\section{Protocol}\label{sec:usd-protocol}

\subsection{The information set on each date}

Fix an observation date $t$. The \emph{analyst's information set} $\mathcal I_t$ is declared to consist of exactly the following, and of nothing else:
\begin{enumerate}[nosep,label=(\alph*)]
\item the quoted predecessor caplet smile on three-month \gls{libor}, at every expiry in the grid of Section~\ref{sec:usd-ident}, as of the close of date $t$, together with the vendor's quoted-or-derived flag for each point;
\item the \gls{libor} forward curve and the \gls{sofr} \gls{ois} discount and forward curves as of the close of date $t$, from which the forward basis $B_0(t)=L_0(t)-F_0(t)-s$ is read;
\item the fallback spread $s$, a constant fixed on 5~March 2021 and taking the value recorded for the three-month tenor in Appendix~\ref{app:data};
\item the realised basis series -- the \gls{libor} fixing minus compounded \gls{sofr} over the matching period minus $s$ -- for every value date up to and including $t-1$;
\item the calendar of scheduled Federal Open Market Committee meeting dates, which is published in advance.
\end{enumerate}

The following are \emph{hidden} and may not enter any quantity computed at date $t$:
\begin{enumerate}[nosep,label=(H\arabic*)]
\item\label{it:H1} every \gls{sofr} option quote, at date $t$ and at every other date, whether from the vendor cube or from the exchange;
\item\label{it:H2} every observation dated later than $t$, of any series, including the realised basis;
\item\label{it:H3} the outcome of any test in this chapter, which is why every threshold is fixed in Chapter~\ref{ch:identifiability} or in this section and none is revised afterwards;
\item\label{it:H4} the identity of the first month $m^*$ in which the vendor marks the successor surface as quoted, when that month defines a window: it is taken from the vendor's own flag as recorded in Appendix~\ref{app:data}, never by inspection of the predictions.
\end{enumerate}

Restriction~\ref{it:H1} is what makes the exercise a prediction rather than a fit, and it is also the second of the two non-negotiable design constraints stated in Section~\ref{sec:ident-validation}: the basis volatility that sets the width of the interval is estimated from the basis record alone.

One exemption to~\ref{it:H1} is made, and one only. Criterion~V4 compares the naive mark with the point estimate obtained \emph{after $v_1$ has been estimated from a single successor quote at the shortest expiry in the grid}, so it cannot be run under~\ref{it:H1} as stated. The exemption: on each observation date the at-the-money successor quote at the shortest expiry is visible, is used to estimate the single coordinate $v_1=\eta\rho_{BF}$ by Step~4$'$ below, and is used for nothing else. Three consequences are recorded rather than argued away. The anchoring pair enters neither V4 test, its error being zero by construction, and the dates on which that quote is missing or derived are counted. Only V4 becomes a one-parameter-anchored comparison; V1, V2, V3, V5 and V6 stay under the full restriction. And since the same $v_1$ is carried to every other expiry, V4 rests on the shared-basis assumption that V6 tests, so the two verdicts are reported together.

Restriction~\ref{it:H2} is enforced mechanically: every trailing statistic is computed on a window ending at $t-1$, which also accommodates the one-business-day lag with which the compounded rate over a period ending at $t-1$ becomes known.

\subsection{The windows}

Five windows are fixed. Each is justified from the cessation calendar and from the data constraints of Appendix~\ref{app:data}, and from nothing else. In particular no window boundary is placed to include or to exclude any episode of market stress; where such an episode falls inside a window, that is a property of the window and is reported in the diagnostics, not a reason to move it.

\paragraph{$W_{\mathrm{pred}}$, the prediction window.} From 5~March 2021, the date on which the spread was fixed and after which $s$ is a constant, to 31~December 2021, the last date on which the predecessor surface is quoted under the ``no new \gls{libor}'' guidance rather than increasingly derived from the successor surface. Before the spread fixing, $s$ was a forecast and the strike shift of Definition~\ref{def:ident-observable} would carry the analyst's own estimate of it; after the end of 2021 the predecessor surface is contaminated in the reverse direction. Both boundaries come from the calendar.

\paragraph{$W_{\mathrm{obs}}$, the observation window.} From $m^*$, the first month in which the vendor marks the successor cap-and-floor surface as quoted rather than derived, to 31~December 2024. The lower endpoint is not yet known and is not chosen here; Appendix~\ref{app:data} records it as an item of the terminal checklist, with the constraint that it is not earlier than November 2021. The upper endpoint is the end of the data range requested.

\paragraph{$W_{\mathrm{ident}}$, the identification-test window.} The intersection $W_{\mathrm{pred}}\cap W_{\mathrm{obs}}$, on which both the predecessor smile is quoted and the successor smile is quoted. It is the strict arm of the test, and it may be short: if $m^*$ falls in November 2021 it is about six weeks of business days. Its length is a finding about the transition, not a defect of the design, and it is reported as such.

\paragraph{$W_{\mathrm{ex}}$, the exchange arm.} The whole of $W_{\mathrm{pred}}$, with the observed successor smile taken from settlement prices of listed options on three-month \gls{sofr} futures rather than from the vendor cube. These rest on a market quoted by construction whose history begins before $m^*$, so the arm extends the comparison backwards at the cost of a shorter expiry reach and an exercise-style conversion. It is reported separately and never pooled with the strict arm.

\paragraph{$W_{\mathrm{cert}}$, the certificate windows.} The fitting sample $\Dfit$ is the business days of calendar 2021, on which the policy is fixed; the calibration record is calendar 2022, separated from $\Dfit$ by a gap of $b$ periods; the deployment window is calendar 2023, separated by a further gap of $b$, extended through 2024 where the data range allows, which increases the number of deployment blocks without changing the threshold.

\paragraph{$W_\eta$, the basis-volatility window.} A trailing window of $252$ business days of the \emph{realised} basis, ending at $t-1$. The quoted basis-swap series is not used for this purpose at any date: quoted spreads compress mechanically toward $s$ in the final year before June 2023, cleared contracts having been converted at the fixed spread, so the last part of that series measures a convergence trade and not a market price of basis risk. The quoted series is used only as a drift check on $B_t$ over the sub-period ending in mid-2022, reported as a diagnostic.

\datanote{Six series carry the windows of this section and all six are described in Appendix~\ref{app:data} and nowhere here: the daily predecessor fixing and the daily successor overnight series, from which the compounded rate over each accrual period is recomputed inside the package; the enacted fallback spread adjustments, of which only the three-month value is used; the daily predecessor cap-and-floor volatility strip with the vendor quoted-or-derived flag on every point; the same strip on the compounded successor rate from $m^*$; and the settlement prices and implied volatilities of listed options on the first eight quarterly successor futures. The date ranges requested are those of the windows above. The one item with no substitute is the predecessor volatility strip.}

\gap{Citation needed: the announcement of 5~March 2021 fixing the spread adjustments and triggering the cessation timetable; the confirmation that the panel ended on 30~June 2023; the ``no new \gls{libor}'' supervisory guidance that bounds $W_{\mathrm{pred}}$ at the end of 2021; and the industry initiative that moved non-linear trading to the successor rate on 8~November 2021, which is the reason $m^*$ cannot be earlier than that month. None of these documents has been obtained and none has a bibliography entry, so the four dates stand as unverified, and what rests on each is stated exactly. The 31~December 2021 and 8~November 2021 dates place window boundaries and nothing else. The 5~March 2021 date additionally carries the constancy of $s$, which is what makes the realised basis $B^{\mathrm{real}}=\text{fixing}-\text{compounded successor}-s$ well defined; $B^{\mathrm{real}}$ is the sole input to $\hat\eta$ and $\hat\eta$ sets the half-width that V1, V3 and V5 test. The mitigation is exact rather than asserted: $\hat\eta$ is a volatility of first differences and so is invariant to the value of $s$, and the date matters only if $s$ was still moving inside $W_\eta$, for which the quoted-basis drift check is the diagnostic. The 30~June 2023 date additionally fixes the record length of about eleven hundred periods used in criterion~R2, and nothing else.}

\subsection{Reporting rules fixed in advance}

\paragraph{Quoted before derived.} A point enters a computation only if the vendor flags it as quoted; where no flag is exposed, the rule of Appendix~\ref{app:data} treats strikes beyond one hundred basis points from the money as unquoted. The discarded count is reported by window and expiry, a rule that discards most of the grid changing what the containment rates mean.

\paragraph{Observation dates and expiries.} The observation dates are the Wednesday marks of each week in the window, or the preceding business day where a Wednesday is a holiday, and the expiry grid is $\{1,2,3,5,7,10\}$ years; both are fixed in Section~\ref{sec:ident-validation} and are repeated here only because they set the denominators below. If the data support only expiries out to two years the grid truncates to $\{1,2\}$, the truncation is reported with the result, the criteria are evaluated on the truncated grid, and no criterion is restated at a new threshold. One consequence is fixed here rather than discovered afterwards. Criterion~V4 passes only if its test over expiries of three years and longer rejects, and item~F5 falsifies the improvement claim only on that same arm; under truncation that arm is empty, so V4 is recorded as \emph{not evaluable} and the improvement claim as \emph{not tested} -- not passed, not failed -- with the pooled test carrying no verdict. The minimum viable export of Appendix~\ref{app:data} reaches two years on the exchange arm at best, so this is the expected outcome there and not a remote contingency. The grid obtained is a row of Table~\ref{tab:usd-ident-summary}.

\paragraph{Order of operations.} Fixed thresholds leave the order of examination free, and that is the last degree of freedom. It is fixed here: strict arm before exchange arm; pooled row before per-expiry rows; the primary cell $(b,C,r,\varepsilon)=(10,1,2,0.05)$ before the rest of the grid; and the falsification rules of Section~\ref{sec:usd-falsify} read before any verdict is written.

\paragraph{Scheduled policy dates.} Forward rates on the successor benchmark move in steps at scheduled monetary-policy dates and diffuse between them, a feature the continuous-forward description of Chapter~\ref{ch:prelim} does not carry \citep{brace2024}. The rule fixed here is that the identification test uses expiries of one year and longer only, so that every accrual window contains several scheduled dates and the step component is absorbed into the fitted volatility on both sides. That the absorption is common to the two sides, and so cancels at the order of the expansion, is an assumption and not a result.

\gap{Not established here: that a jump component at scheduled policy dates, present in both rates with the same timing, leaves the contamination map of Chapter~\ref{ch:identifiability} unchanged at first order in the basis scale. That chapter records that a jumping \emph{basis} would add a third-cumulant term proportional to neither contamination direction; the present statement is the weaker one about a jump in the \emph{rate} common to both benchmarks. Neither is proved.}

\paragraph{Exchange-arm conversion.} Settlement prices are converted to normal implied volatilities on the compounded rate with the contract's own settlement convention and, if the options are American-style, with an early-exercise adjustment whose size is reported for every point rather than assumed negligible.

\gap{Citation needed: the contract specification of the listed three-month successor-rate futures and of the options on them, for the settlement formula, the exercise style and the listing dates. No primary document has been obtained; until one is, the exchange arm is specified but its conversion step is not fixed.}

\section{The identification test}\label{sec:usd-ident}

\subsection{The quantities, defined precisely}

Fix an observation date $t\in W_{\mathrm{ident}}$ and an expiry $T$ in the grid. Write $\Delta$ for the three-month accrual length, $\tau^*=T+\Delta/3$ for the clock of Chapter~\ref{ch:prelim} and $q=\sqrt{\tau^*/(T+\Delta)}$ for the quoting factor of Chapter~\ref{ch:identifiability}.

\paragraph{Step 1: the fitted predecessor triple.} Let $\mathcal K_t$ be the set of strikes at expiry $T$ that the vendor flags as quoted on date $t$, and let $\sigma^{N,\mathrm{obs}}_L(K)$ be the quoted normal implied volatility at each. The fitted triple is
\[
(a,r,n)(t,T):=\arg\min_{(a',r',n')}\ \sum_{K\in\mathcal K_t}w_K\Big(\sigma^{N,\mathrm{obs}}_L(K)-\sigma^{N,\mathrm{SABR}}(K;a',r',n',T)\Big)^2 ,
\]
where $\sigma^{N,\mathrm{SABR}}$ is the normal \gls{sabr} implied volatility of Chapter~\ref{ch:prelim} at expiry $T$ with $\beta=0$, and the weights are fixed in advance as $w_K\equiv1$, that is, an unweighted least-squares fit in volatility units over the quoted strikes. The minimisation is over $a'>0$, $|r'|<1$, $n'>0$. The fit residual in basis points is recorded for every $(t,T)$ and is reported as a diagnostic; it is not an acceptance criterion, and Chapter~\ref{ch:identifiability} explains why a small residual is not evidence that the triple is the successor's triple.

\paragraph{Step 2: the strike shift and the time change.} The predicted successor smile is the map
\[
\Phi_0(K'):=\sigma^{N,\mathrm{SABR}}\big(K'; \alpha_0,\rho,\nu,\tau^*\big)\cdot q ,\qquad K'=K-s-B_0(t),
\]
with $(\alpha_0,\rho,\nu)$ the observable-based values of Section~\ref{sec:ident-validation}, namely $\alpha_0=\sqrt{a^2-\hat\eta^2}$, $\nu=na^2/\alpha_0^2$ and $\rho$ from \eqref{eq:ident-inverse} at $p=0$. The factor $q$ is the conversion from the clock to the quoted convention of Chapter~\ref{ch:prelim}. The \emph{naive mark} is the same construction with $\hat\eta$ set to zero, which is the market convention: fit, shift the strike, apply the decay, quote.

The successor fixing is paid with the two-day backward shift of the benchmark administrator's rule book recorded in Appendix~\ref{app:data}, so the payment date is not $T+\Delta$. The shift is absorbed rather than modelled, on the bound of Chapter~\ref{ch:identifiability}: under $\Q^{T+\Delta+\delta_p}$ the forward acquires a drift equal to minus its covariation with the logarithm of the numeraire ratio, which is $\delta_p$ times the stub forward at first order, so the displacement of the expected forward is at most $\alpha_0\sigma_{\mathrm{stub}}\delta_p(T+\Delta)\le0.005$ bp at $\alpha_0=\sigma_{\mathrm{stub}}=60$ bp, $\delta_p\le4$ calendar days and $T+\Delta=1.25$, against a quote precision of $0.01$ bp.

\datanote{The stub volatility $\sigma_{\mathrm{stub}}$ is estimated nowhere in this thesis, so the bound above is a bound on an unmeasured quantity at an assumed level of $60$ bp, that of the worked example of Chapter~\ref{ch:identifiability}. It is linear in $\sigma_{\mathrm{stub}}$, so a doubling of the assumed level brings the displacement to the quote precision itself and the dismissal is safe only to about that factor.}

\paragraph{Step 3: the basis volatility.} Let $B^{\mathrm{real}}_u$ denote the realised basis at value date $u$. The estimator is the trailing realised volatility
\[
\hat\eta(t):=\sqrt{\frac{252}{N_w}\sum_{i=1}^{N_w}\big(B^{\mathrm{real}}_{u_i}-B^{\mathrm{real}}_{u_{i-1}}\big)^2},\qquad N_w=252,
\]
over the window $W_\eta$ ending at $t-1$, with no mean removed, the basis being a martingale under the pricing measure. Section~\ref{sec:ident-validation} transmits its relative standard error one-for-one to the half-width; the band is reported beside every containment rate and is folded into no criterion.

\paragraph{Step 4: the identified interval.} The interval $[\Sigma^+(t,T),\Sigma^-(t,T)]$ is the one fixed in criterion~V1 of Section~\ref{sec:ident-validation}, with $\hat\eta(t)$ in place of $\eta$ and each extreme model carried through its own Hagan bracket.

What must be fixed here is admissibility, because $\hat\eta$ is a free estimate and Theorem~\ref{thm:ident-two} has hypotheses it can violate on any date: $a>\hat\eta\sqrt2$, without which $\alpha_0(p)=\sqrt{a^2-\hat\eta^2(1-p^2)}-p\hat\eta$ is not real along the homotopy and $\Sigma^+$ loses its meaning, and the side condition $|r|<\sqrt{1-\hat\eta^2/a^2}$, under which the intermediate models carry a correlation below one in modulus. They also control the ordering of the endpoints, since the bracket of the closed-form half-width $\hat\eta q[1-(2-3r^2)n^2\tau^*a^2/(24(a^2-\hat\eta^2))]$ in \eqref{eq:ident-quotedhw} falls as $\hat\eta\to a/\sqrt2$, so at large $\hat\eta$ the asserted $\Sigma^+\le\Sigma^-$ can fail. A pair $(t,T)$ is therefore excluded when $\hat\eta(t)\ge a(t,T)/\sqrt2$ or $|r(t,T)|\ge\sqrt{1-\hat\eta(t)^2/a(t,T)^2}$, with the count reported as a column of Table~\ref{tab:usd-ident-summary}. This is the upper counterpart of the lower exclusion $\hat\eta<1$ bp fixed for V5: an interval reported where its theorem does not apply is not evidence about the theorem.

\paragraph{Step 4$'$: the $v_1$-anchored point estimate.} Criterion~V4 needs a point in the interval, not the interval. Let $T_{\min}$ be the shortest expiry and $\Sigma^{\mathrm{anc}}(t)$ the anchoring quote released in Section~\ref{sec:usd-protocol}. Along the homotopy \eqref{eq:ident-homotopy} the map $w\mapsto(a-w)c^*q$ is continuous and strictly decreasing from $\Sigma^-$ to $\Sigma^+$ by \eqref{eq:ident-alpharange}, so $\hat v_1(t)$ is the unique solution of $(a(t,T_{\min})-\hat v_1)c^*q=\Sigma^{\mathrm{anc}}(t)$, found by bisection and set to the nearer endpoint, with a flag, when the quote lies outside the interval. Since $v_1=\eta\rho_{BF}$ is shared across expiries, the $v_1$-anchored estimate at every other expiry is $\Sigma^{v_1}(t,T):=(a(t,T)-\hat v_1(t))c^*q$, the same construction at the same $\hat v_1$, and V4 is the pre-registered sign test of $|\Sigma^{\mathrm{obs}}-\Sigma^{v_1}|$ against $|\Sigma^{\mathrm{obs}}-\Phi_0(F_0)|_{\hat\eta=0}$.

The naked midpoint $\tfrac12(\Sigma^++\Sigma^-)$ is not that estimate and may not stand in for it. Each endpoint carries its own bracket, so the midpoint is $q[a+(2-3r^2)n^2a^3\tau^*/(24(a^2-\hat\eta^2))]$ against the naive mark $q[a+(2-3r^2)n^2a\tau^*/24]$, a relative difference of exactly $\tfrac1{24}(2-3r^2)n^2\tau^*\hat\eta^2/(a^2-\hat\eta^2)$, which is $0.0865\%$ at the parameters of Section~\ref{sec:ident-example}, or $0.056$ bp on a $64.6$ bp level. A sign test between two marks separated by five hundredths of a basis point, on a market quoted to one hundredth, decides on rounding; the midpoint is kept as a diagnostic row only.

\paragraph{Step 5: the observation.} Let $\Sigma^{\mathrm{obs}}(t,T)$ be the quoted successor at-the-money normal volatility at the same date and expiry, taken from the vendor cube in the strict arm and from exchange settlements in the exchange arm. The quoted successor skew $\varsigma^{\mathrm{obs}}(t,T)$ is the slope at the money of the ordinary least-squares line fitted to the quoted successor normal volatilities over the strikes within fifty basis points of the money, a definition fixed here in advance because the criterion of Section~\ref{sec:ident-validation} refers to ``the quoted successor skew'' without fixing an estimator. Its predicted counterpart is $\Phi_0'(F_0)$, the analytic slope of the predicted smile at the money.

\subsection{The statistics}

Three statistics carry the test. Write $\mathcal P$ for the set of observation-date-and-expiry pairs in the window, $|\mathcal P|=M_{\mathcal P}$.

\begin{definition}[Containment rate, normalised excess, relative half-width]\label{def:usd-stats}
For $(t,T)\in\mathcal P$ set
\[
C_{t,T}:=\ind{\Sigma^{\mathrm{obs}}(t,T)\in[\Sigma^+(t,T),\Sigma^-(t,T)]},\qquad
w_{t,T}:=\tfrac12\big(\Sigma^-(t,T)-\Sigma^+(t,T)\big),
\]
and define the \emph{signed normalised excess}
\[
D_{t,T}:=\frac{\big(\Sigma^{\mathrm{obs}}-\Sigma^-\big)^+-\big(\Sigma^+-\Sigma^{\mathrm{obs}}\big)^+}{w_{t,T}} ,
\]
which is zero exactly when the pair is contained, positive when the market quoted above the interval and negative when below, and measured in units of the half-width. The \emph{containment rate} is $\bar C:=M_{\mathcal P}^{-1}\sum_{\mathcal P}C_{t,T}$, computed overall and separately by expiry. The \emph{relative half-width} is $w_{t,T}/\Phi_0(F_0)(t,T)$.
\end{definition}

The normalised excess distinguishes a test that fails narrowly from one that fails grossly, which is why $\bar C$ is never reported alone: a containment rate of $0.85$ with $|D|$ below $0.1$ on the failures is a different finding from the same rate with $|D|$ above $2$, and only the second is evidence against the model class.

\paragraph{Dependence correction.} The pairs are strongly dependent, consecutive Wednesday marks sharing almost all their information and the expiries at one date sharing the fitted triple. Every rate and test statistic below therefore carries a moving-block bootstrap interval, resampling whole observation dates so that all expiries of a date move together, with block length $b^*$ evaluated by Theorem~\ref{thm:block-b} at the record length in dates and $r=2$, $R=2000$ replications and the seed of Table~\ref{tab:appC-seeds}. The interval is a precision statement beside each criterion and does not replace it: the criteria of Section~\ref{sec:ident-validation} are raw proportions and fixed-level tests, evaluated exactly as stated there.

\subsection{Acceptance criteria, inherited}

The six criteria V1 to V6 are fixed in Section~\ref{sec:ident-validation} of Chapter~\ref{ch:identifiability}, before any data were pulled, and this chapter changes none of them. They are reported in full, each with its verdict, and the count of criteria passed is reported as a single line so that a reader cannot assemble a favourable subset. Three points of execution need fixing here, and all three are matters of computation rather than of threshold.

First, V4 names its estimator only by description, and Step~4$'$ supplies it as $\Sigma^{v_1}$, with the exemption to~\ref{it:H1} it requires declared in Section~\ref{sec:usd-protocol}. V4 is also a paired sign test over serially dependent pairs, so its nominal five per cent level is not its true level: the criterion is evaluated exactly as written on $\Sigma^{v_1}$, the block-bootstrap $p$-value for the same null is reported beside it, and disagreement is recorded as inconclusive. Both pre-registered tests are run and only the long-expiry one decides the criterion.

Second, V5's implied $\rho_{BF}$, the realised deviation divided by $-\hat\eta qc^*$, is undefined at $\hat\eta=0$ and unstable at small $\hat\eta$: a pair is excluded from V5 when $\hat\eta(t)$ falls below one basis point, with the count reported, and the implied value is truncated to $[-1,1]$ before any correlation is computed, values outside that range not being correlations, with the count and magnitude of truncations reported. A high truncation count is itself evidence, and Section~\ref{sec:usd-falsify} says what of. Truncation is winsorisation, and the Pearson correlation of a winsorised series is not the one Chapter~\ref{ch:identifiability} pre-registered, so the untruncated correlation is reported in the adjacent column and disagreement between the two is inconclusive.

Third, the pairs excluded by each of the three rules -- derived or unflagged points, $\hat\eta$ below one basis point, and the admissibility rule of Step~4 -- are counted before any rate, all three moving the denominators of this section.

\subsection{Tables and figures the test produces}

\begin{table}[htbp]
\centering\small
\caption{Specification of the identification-test summary table, one row per expiry and a pooled row, produced separately for the strict and exchange arms. It serves criteria V1, V2 and V3 and reports the precision of each. No entry is filled here.}
\label{tab:usd-ident-summary}
\begin{tabular}{@{}ll@{}}
\toprule
Column & Content \\
\midrule
Expiry & years; pooled row last \\
Grid obtained & expiries actually supported by the export, and whether the grid truncated \\
$M_{\mathcal P}$ & number of observation-date pairs after the quoted-only rule \\
Discarded & points discarded as derived or unflagged, per cent of grid \\
Inadmissible & pairs excluded by $\hat\eta\ge a/\sqrt2$ or $|r|\ge\sqrt{1-\hat\eta^2/a^2}$ (Step~4) \\
Fit residual & median across pairs of the root-mean-square fit residual, bp \\
$\bar C$ (V1) & containment rate of $\Sigma^{\mathrm{obs}}$ in $[\Sigma^+,\Sigma^-]$ \\
$\bar C$ interval & block-bootstrap $95\%$ interval for $\bar C$ \\
V1 verdict & pass if $\bar C\ge0.90$ \\
$\bar C_{\varsigma}$ (V2) & containment rate of the quoted skew in its interval \\
V2 verdict & pass if $\bar C_\varsigma\ge0.90$ \\
Median $w$ & median half-width, bp \\
Median $w/\Phi_0$ & median relative half-width, per cent \\
V3 verdict & informative if median $w/\Phi_0\le0.25$ \\
$\hat\eta$ band & median $\hat\eta$ and its relative standard error, bp \\
\bottomrule
\end{tabular}
\end{table}

\begin{table}[htbp]
\centering\small
\caption{Specification of the failure-diagnosis table, expiries crossed with calendar quarters. It distinguishes narrow from gross and clustered from dispersed failure, and is the evidence base for the attribution rules of Section~\ref{sec:usd-falsify}.}
\label{tab:usd-ident-excess}
\begin{tabular}{@{}ll@{}}
\toprule
Column & Content \\
\midrule
Cell & expiry $\times$ calendar quarter \\
$M_{\mathcal P}$ & pairs in the cell \\
$\bar C$ & containment rate in the cell \\
$D$ quantiles & $5$th, $25$th, $50$th, $75$th, $95$th percentiles of the signed normalised excess \\
Failures above & share of failures above the interval, $\#\{D>0\}/\#\{D\ne0\}$ \\
Max $|D|$ & largest normalised excess in the cell \\
Naive error & median $|\Sigma^{\mathrm{obs}}-\Phi_0(F_0)|_{\hat\eta=0}$, bp \\
V4 point error & median $|\Sigma^{\mathrm{obs}}-\Sigma^{v_1}|$, bp; the estimator criterion~V4 names \\
Midpoint error & median $|\Sigma^{\mathrm{obs}}-\tfrac12(\Sigma^++\Sigma^-)|$, bp; diagnostic only \\
Implied $\hat\eta$ ratio & median and interquartile range of $\eta^{\mathrm{imp}}/\hat\eta$, where $\eta^{\mathrm{imp}}$ is the basis volatility that would have made the pair exactly contained \\
\bottomrule
\end{tabular}
\end{table}

\begin{table}[htbp]
\centering\small
\caption{Specification of the nuisance-parameter table, one row per expiry. It serves criteria V4, V5 and V6, and it is the only place where a statement about the pricing measure is compared with a statement about the physical measure.}
\label{tab:usd-ident-rhobf}
\begin{tabular}{@{}ll@{}}
\toprule
Column & Content \\
\midrule
Expiry & years; then the two pooled rows of criterion~V4 \\
Sign-test $p$ (V4) & one-sided paired sign test, $\Sigma^{v_1}$ against the naive mark \\
Bootstrap $p$ (V4) & block-bootstrap $p$-value for the same null \\
V4 pooled rows & the whole grid, and expiries of three years and longer; the second decides V4 \\
Ties discarded & pairs tied to $0.01$ bp; and the anchoring pairs, which enter neither test \\
Implied $\rho_{BF}$ & median and interquartile range across dates \\
Truncated & pairs with implied value outside $[-1,1]$, and median excess \\
Realised $\rho_{BF}$ & median trailing correlation of the basis with the compounded forward \\
Pearson $\rho$ & correlation of the two series across dates, truncated and untruncated \\
Bootstrap $p$ (V5) & moving-block bootstrap $p$-value against zero, for each of the two \\
Sign agreement & whether implied and realised agree in sign \\
IQR (V6) & interquartile range of the implied series \\
Max median gap (V6) & maximum absolute difference between the per-expiry medians; \emph{stable} needs this below $0.2$ as well as the IQR \\
Slope on $B_0$ (V6) & ordinary least-squares slope and its bootstrap $p$-value \\
V6 verdict & stable / moves with the level / neither \\
\bottomrule
\end{tabular}
\end{table}

\figplan{Four panels, one per expiry in $\{1,2,5,10\}$ years, on one observation date fixed in advance as the first Wednesday of the second month of $W_{\mathrm{ident}}$. Horizontal axis: strike offset from the successor forward, basis points, over $[-150,150]$. Vertical axis: normal implied volatility, basis points. Series: the predicted smile $\Phi_0$ solid; the two endpoint smiles dashed with the region between them shaded; the naive mark dotted; the quoted successor smile as filled markers, open where the vendor flags a point as derived; a vertical rule at the money.}{The predicted successor smile, its identified band and the observed successor smile at one date. It tests by eye what criterion~V1 measures over the whole window, and shows the band's width against the distance between the naive mark and the observation, which is what criterion~V4 tests.\label{fig:usd-band}}

\figplan{Six panels, one per expiry, for the strict arm; the exchange arm goes in a second figure of identical construction on the wider axis of $W_{\mathrm{pred}}$ but of \emph{two} panels only, at one and two years, the listed options reaching two years on the first eight quarterly contracts and the policy-date rule of Section~\ref{sec:usd-protocol} removing everything below one year -- the emptiness of the rest is itself the finding. Horizontal axis: observation date. Vertical axis: at-the-money normal implied volatility, in basis points. Series: the shaded band between $\Sigma^+$ and $\Sigma^-$; the predicted midpoint solid; the naive mark dotted; the quoted successor at-the-money volatility solid with markers; breaches marked; and a rug marking dates on which more than half the grid was discarded by the quoted-only rule.}{Containment through time, by expiry. The figure tests whether failures of criterion~V1 are dispersed or clustered, which is what distinguishes a mis-specified model class from a stale basis-volatility estimate; the rug shows where the quoted-only rule has thinned the evidence, so that a run of containment on a thin grid is not read as one.\label{fig:usd-containment}}

\figplan{Two panels. Left: empirical distribution functions of the signed normalised excess $D$, one curve per expiry, horizontal axis $D$ over $[-3,3]$, with vertical rules at $\pm1$ marking the band edges in normalised units. Right: $D$ scattered against the relative half-width $w/\Phi_0$ in per cent, points coloured by expiry, with a vertical rule at the criterion-V3 threshold of twenty-five per cent.}{Failure magnitude and its relation to interval width. The left panel tests whether failures of criterion~V1 are narrow or gross; the right tests the joint hypothesis that the interval is at once informative and containing, which is criteria V1 and V3 read together and the only reading under which the identification result is operationally useful.\label{fig:usd-excess}}

\figplan{Two panels on a shared date axis over $W_{\mathrm{pred}}$ and $W_{\mathrm{obs}}$. Upper, in basis points: the trailing $\hat\eta(t)$ with a shaded band of one relative standard error $\hat\eta(t)/\sqrt{2N_w}$, and the resulting median half-width across expiries on a secondary axis. Lower, in units of correlation over $[-1,1]$: the implied $\rho_{BF}$ series of criterion~V5, one line per expiry, the trailing realised correlation of the basis with the compounded forward overlaid as a heavy line, and markers where the implied value has been truncated.}{The one estimated input and the one testable nuisance parameter. The upper panel tests how much of the interval's width is driven by movement in the basis-volatility estimate rather than by the fitted predecessor smile; the lower is criteria~V5 and~V6 in one picture, agreement in sign and level between the option-implied and historically estimated basis-rate correlation and the stability of the implied series. The correlation window is the $252$-business-day trailing window of $\hat\eta$.\label{fig:usd-eta}}

\datanote{Every quantity in this section is computed from four inputs and no others: the quoted predecessor smile with its quoted-or-derived flags, the two curves, the constant spread, and the realised basis series. The realised basis is computed inside the reproducibility package from the daily overnight fixings published by the central bank, compounded over each accrual period with the conventions recorded in Appendix~\ref{app:data}, minus the predecessor fixing from the terminal export, minus the constant spread; only the predecessor fixing is not redistributable. The fit of Step~1 is a three-parameter least-squares problem solved by the module described in Appendix~\ref{app:data}; the starting values, the tolerance and the treatment of non-convergence are recorded there, and any pair for which the fit does not converge is excluded and counted.}

\section{The certificate test}\label{sec:usd-cert}

\subsection{The policy, and why it is the market convention}

Theorem~\ref{thm:cov} requires only that the valuation-and-hedging policy be fixed on the fitting sample; it makes no assumption that the model underlying the policy is correct. That licence is used here in the strongest way available: the policy is the market convention, exactly as a desk would implement it, with no use of anything proved in this thesis.

\begin{definition}[The certified policy]\label{def:usd-policy}
The policy $\pi$ is: on each date, fit a normal \gls{sabr} triple to the quoted predecessor caplet smile by the unweighted least-squares fit of Step~1 of Section~\ref{sec:usd-ident}; shift the strike by the spread plus the forward basis; apply the linear volatility decay with the $\Delta/3$ effective maturity; value each successor caplet in the book at the resulting normal volatility; and hedge with a delta in the successor forward computed from the same fitted triple, rebalanced at each daily close. The parameters of the decay convention and the hedge ratio are functions of $\Dfit$ only through the fit on the last date of $\Dfit$, and are frozen thereafter for the purposes of the certificate.
\end{definition}

The policy is delta-only, and the choice is forced. A vega hedge is a position in options, and over $\Dfit$ and the early calibration window there is no quoted successor option market -- the chapter's own premise and the reason $m^*$ exists -- so a vega leg would be untradeable on exactly the window on which the policy is fixed, and a certificate on a policy that could not have been run is a certificate on nothing. The vega-hedged variant is therefore a \emph{second} policy, certified separately over the same windows, with its vega leg executed in the listed option on the three-month successor future, which trades from before $m^*$ at the cost of an expiry reach shorter than the book's. The two are reported side by side, neither as the other's correction.

Two features of the choice matter. It is not the thesis's own policy -- no identified interval, midpoint or reserve enters it -- and it is not a good policy in any sense the thesis defends, being the policy whose hedging error the thesis exists to bound. The certificate is indifferent to which fixed policy is certified, which is exactly why the market convention is the informative choice.

\paragraph{The book.} The hedged portfolio is a fixed notional of successor caplets at each expiry in the grid, struck at the money on the first date of $\Dfit$ and held to the end of the deployment window, with no re-striking and no roll.

The per-period hedging error has two legs and is their difference. The \emph{risk-theoretical} leg $\Delta V^{\mathrm{rt}}_t$ is the daily change in the value of the portfolio and its hedges with every position revalued at the policy's own volatilities, which is the change the policy predicts and is available on every date because the policy manufactures its own marks. The \emph{hypothetical} leg $\Delta V^{\mathrm{hyp}}_t$ is the daily change in the value of the same positions revalued at the observed successor marks. Then $e_t:=\Delta V^{\mathrm{rt}}_t-\Delta V^{\mathrm{hyp}}_t$, with the sign convention of Chapter~\ref{ch:certification} that positive values are losses. Marking both legs at the policy's own volatilities would make $e_t$ the discretisation residual of the policy's own model and would measure nothing about the market; these two legs are exactly those the regulatory attribution test of Section~\ref{sec:usd-methods} compares, which is why that test is the natural comparator and not an analogy. The hypothetical leg needs observed successor prices, so $e_t$ exists only from $m^*$, which is why the deployment window lies entirely after $m^*$. Before $m^*$ only the risk-theoretical leg is computable; it is kept as a diagnostic and enters no score, threshold or criterion. The block score is the cumulative block shortfall, $\phi(p,x)=\sum_i x_i$, which is the default of Definition~\ref{def:score} and is the quantity a reserve is held against; the absolute cumulative error is reported as a check, as Chapter~\ref{ch:certification} recommends.

\datanote{The book is synthetic and declared here rather than taken from any institution: equal vega at each expiry in $\{1,2,3,5,7,10\}$ years on the first date of $\Dfit$, long only. Nothing in the certificate depends on the composition, which enters only through the law of the score; it is stated so that the reported thresholds are reproducible in currency units. The risk-theoretical leg is available on every date because the policy manufactures its own marks, which is not a use of hidden data; the hypothetical leg needs the observed successor surface and so exists only from $m^*$, and the calibration window is reported with the count of dates on which it rests on a derived surface.}

\gap{Not established here: that a calibration record whose hypothetical leg rests in part on a vendor-derived successor surface satisfies the common-law requirement of Assumption~\ref{ass:mix}. If the count of such dates exceeds one quarter of the calibration dates -- a trigger fixed here in advance -- the calibration window is shortened to begin at $m^*$ and the loss of blocks is reported. That is the only window revision this chapter permits, and it is triggered by a data property, never by a result.}

\subsection{Blocks, and a block length the record cannot supply}

The layout is that of Chapter~\ref{ch:certification}: the fitting sample, a gap of $b$ periods, $n=\lfloor N/(2b)\rfloor$ calibration blocks alternating with gaps of $b$, a trailing gap, then the deployment blocks. With the calibration window a calendar year of daily observations, $N$ is about $250$.

The feasibility constraint on $n$ bites at once. At $\alpha=0.1$ and $N=250$ the block length $b=13$ gives $n=9$ and $k=9$, so the threshold is the largest of nine calibration scores and $k/(n+1)$ is exactly $0.9$; at $b=14$ the threshold is infinite and the certificate is vacuous, so the admissible range is $b\le13$. Against that, the block length minimising the expected calibration-conditional shortfall of Chapter~\ref{ch:sharpness}, $b^*=(2(r+1))^{2/(2r+3)}N^{3/(2r+3)}$, is about $17.8$ at $r=2$, and keeping the discarded constants moves it further up. The record is too short to carry the block length the certificate bound prefers.

This is reported as a finding and not repaired. The certificate is evaluated at $b\in\{1,2,5,10,13\}$, giving $n=125,62,25,12,9$ and $k/(n+1)=0.905,0.905,0.923,0.923,0.900$, with every deficit term reported at each. The gap between $1-\alpha$ and $k/(n+1)$ is a visible finite-$n$ effect and is tabulated beside the nominal level. It is not one displaced target but three: $1-k/(n+1)$ is $0.0952$ at $b=1$ and $b=2$, $0.0769$ at $b=5$ and $b=10$, and exactly $0.1000$ at $b=13$, so the ceiling in $k$ moves the effective target to between seven and a half and ten per cent depending on $b$, and a realised breach frequency must be read against the value in the column beside it, not against the nominal ten per cent.

\paragraph{The deployment blocks.} Theorem~\ref{thm:cov} certifies one deployment block, and every frequency below is a frequency over $M$ of them, so the layout is fixed here: $M=\lfloor N_{\mathrm{dep}}/(2b)\rfloor$ blocks of length $b$ alternating with gaps of length $b$, the first separated from the calibration record by the trailing gap. On 2023 alone $N_{\mathrm{dep}}$ is about $250$ and $M=12$ at $b=10$; extended through 2024 it is about $500$ and $M=25$. The per-block statement is Theorem~\ref{thm:cov}(i) applied $M$ times and does not degrade in $M$, so what C1 bounds is the \emph{expected} breach frequency $\alpha+\max_i\dK(P,Q^{(i)})+(n+1)\beta(b)$ of Remark~\ref{rem:sequence}, not the realised one.

Two consequences are stated rather than assumed away. First, C1's one-sided binomial test treats the $M$ indicators as independent, which the certificate does not supply: by Remark~\ref{rem:sequence} the coupling that would make the $M$ deployment copies independent of the threshold and of each other costs $(n+M)\beta(b)$, which at $b=10$, $C=1$, $r=2$, $M=25$ is $0.37$ against the $0.13$ that C1 charges. That is why the block-bootstrap $p$-value stands beside the binomial one as a co-equal statistic and disagreement is inconclusive.

Second, an $M$ of this size decides what the test can see, and the arithmetic needs no data. At $M=12$ against the primary cell's bound of $0.28$, the test rejects at five per cent only on seven or more breaches out of twelve, a realised frequency of $0.583$, and reaches power one half only against a true breach probability of $0.54$; at $M=25$ it rejects on twelve or more, a frequency of $0.48$, with power one half at $0.46$. Against the nominal $\alpha=0.1$ the two thresholds are $0.333$ and $0.240$. The test detects catastrophic failure and nothing milder, and C1 is reported with these numbers attached so that a pass reads as the absence of catastrophe rather than as confirmation.

\begin{remark}[The block length is chosen by feasibility, not by optimality]\label{rem:usd-b}
Chapter~\ref{ch:sharpness} is explicit that $b^*$ minimises an upper bound on the certificate's own shortfall, no lower bound being proved for the blocked route, so $b^*$ is optimal for the bound and not for the problem. Here the question does not arise: the admissible set does not contain $b^*$ and the choice is made by feasibility of a finite threshold at the stated level. The sweep over $b$ is reported in full so that the sensitivity of every conclusion to it is visible.
\end{remark}

\subsection{The three deficit terms, reported separately}

Theorem~\ref{thm:cov}(i) certifies the level $1-\alpha-\dK(P,Q)-(n+1)\beta(b)$ and part~(ii) the conditional level $b_{n,k}(\delta)-\dK(P,Q)-\beta(b)/\delta'$ at confidence $1-\delta-\delta'-n\beta(b)$. Three quantities are subtracted, they have different statuses, and the chapter reports each on its own line.

\paragraph{The finite-sample term.} The gap $1-\alpha-b_{n,k}(\delta)$ is computed exactly, by bisection on the binomial tail as in Algorithm~\ref{alg:cert}, at $\delta=0.05$; nothing is assumed for it, the Gaussian approximation to the Beta quantile being poor at these block counts and used nowhere. This is the only one of the three that shrinks with the record.

\paragraph{The coupling term.} The terms $(n+1)\beta(b)$ and $\beta(b)/\delta'$ are computed from an assumed class $\beta(b)\le Cb^{-r}$ at $\delta'=0.05$. The coefficient is not estimable distribution-free, as Chapter~\ref{ch:certification} establishes by construction, so $(C,r)$ is an input and never an output; the pre-registered grid is $C\in\{0.5,1,2\}$, $r\in\{1.5,2,3\}$, and levels are reported over the whole grid rather than at a preferred point.

\paragraph{The drift term.} The Kolmogorov distance $\dK(P,Q)$ is charged in two ways, both reported: as a stress input at $\varepsilon\in\{0,0.02,0.05,0.10\}$, the first being the no-drift reference and not a claim; and as a plug-in by Corollary~\ref{cor:plugin} from the deployment scores and an assumed density bound $L\in\{L_0/2,L_0,2L_0\}$, with $L_0$ the reciprocal of the interquartile range of the calibration scores, declared here in advance. Chapter~\ref{ch:certification} certifies the calibration side of that plug-in and not the deployment side, and the table says so on its face. Its $n^{-1/4}$ rate follows from the square root in Proposition~\ref{prop:w1}, so at $n$ of a few tens the plug-in is coarse: a reason to report the stress form beside it, not to omit either.

\subsection{Acceptance criteria, stated in advance}

\begin{enumerate}[label=\textbf{C\arabic*.},leftmargin=2.2\parindent]
\item \textbf{Marginal coverage.} Over the $M$ deployment blocks, the realised breach frequency $\hat\alpha:=M^{-1}\sum_{i}\ind{S_0^{(i)}>\hat q}$ is consistent, by the test below, with the bound $\alpha+\hat\varepsilon+(n+1)\bar\beta$ on its expectation. The criterion is evaluated at $b=10$ as the primary case, at $C=1$, $r=2$, $\varepsilon=0.05$, all fixed here, and at every other cell of the grid as a sensitivity. It is passed if the one-sided binomial test of $\hat\alpha$ against that bound, at the five per cent level, does not reject in the direction of under-coverage; the block-bootstrap $p$-value for the same null is reported beside the binomial one, and where they disagree the finding is inconclusive.

C1 is informative only where its bound is below one. Over the pre-registered grid of $180$ cells below, $\alpha+\hat\varepsilon+(n+1)\bar\beta\ge1$ in $93$ of them: all $36$ at $b=1$, all $36$ at $b=2$, $20$ of $36$ at $b=5$, one at $b=10$, none at $b=13$. There the criterion is passed by a breach on every deployment block, that is by a threshold of $-\infty$, and a criterion satisfied by the worst possible outcome is not a criterion. Such cells are therefore recorded as \emph{not informative} rather than \emph{pass}, with the count printed in the first column-block of Table~\ref{tab:usd-cert}. At the primary cell the bound is $0.28$ and C1 is live.
\item \textbf{Non-vacuity.} The certified marginal level $1-\alpha-\hat\varepsilon-(n+1)\bar\beta$ must be at least $0.5$ in the primary case; a stated level below a coin toss is recorded as vacuous at that block length and mixing class, and the smallest $b$ clearing $0.5$ is reported. The criterion can fail: at $b=1$, $C=1$, $r=2$ the coupling term alone is $(n+1)Cb^{-r}=126$, so the certificate is vacuous at the shortest block length before any drift is charged. Since $\beta(b)\le1$ always, the assumed class delivers nothing beyond the trivial bound at $b=1$ and $126$ is what $126$ copies of a trivial bound sum to. For the certified level merely to be positive there, at $N=250$ and no drift charge, one needs $\beta(1)\le0.9/126=0.0071$, and to clear one half, $\beta(1)\le0.0032$: a daily state process all but independent at lag one, which is the hypothesis this chapter exists in order not to assume.
\item \textbf{Two-sided error.} The realised breach frequency is also compared with the over-coverage bound of Proposition~\ref{prop:twosided}(a): $\hat\alpha$ below $\alpha-1/(n+1)-\hat\varepsilon-(n+1)\bar\beta$ is a failure of the two-sided statement and is reported as such. Systematic over-coverage is a finding about the threshold's conservatism and is the quantity Chapter~\ref{ch:sharpness} optimises.

C3 carries the mirror image of C1's defect, and worse. Its bound is positive in only $2$ of the same $180$ cells -- $b=10$, $r=3$, $\varepsilon=0$, at $C=0.5$ and $C=1$, where it is $0.0166$ and $0.0101$ -- and negative everywhere else, so $\hat\alpha\ge0$ cannot fall below it and C3 cannot fail. Even there, failure needs $\hat\alpha$ below about $0.017$ with $M$ between twelve and twenty-five, that is exactly zero breaches. At the primary cell the bound is $-0.157$ and the two-sided companion has nothing to test on this record. The primary cell is not moved on that account, since moving it would be choosing a cell for the outcome it can produce; instead C3 is recorded as \emph{not informative} wherever its bound is negative, the count $2/180$ is printed beside it, and the two live cells are named in the table.
\item \textbf{Conditional coverage.} With the calibration sample at its single realised value the deployment breach frequency is the natural estimate of the calibration-conditional coverage, and it is compared with the certified conditional level $b_{n,k}(0.05)-\hat\varepsilon-\bar\beta/0.05$. A single calibration draw gives one realisation of an event holding with probability $1-\delta-\delta'-n\bar\beta$, so a single failure is not evidence against the theorem and a single success not evidence for it; what is measured is whether the certificate was conservative on this record. Being a high-probability statement, C4 therefore carries a recording rule and never a pass or a fail. The outcome is \emph{consistent} when $1-\hat\alpha\ge b_{n,k}(0.05)-\hat\varepsilon-\bar\beta/0.05$, \emph{conservative} when it exceeds that level by more than the width of the exact binomial interval for $\hat\alpha$ at $M$ blocks, and otherwise \emph{below the certified conditional level}, with the caveat above attached. The numbers are fixed by the layout and not by the data: $b_{12,12}(0.05)=0.05^{1/12}=0.7791$ and $b_{9,9}(0.05)=0.05^{1/9}=0.7169$, while $\bar\beta/\delta'=20\bar\beta$ is $0.20$ at $b=10$ and $0.118$ at $b=13$ under $C=1$, $r=2$, so the certified conditional level is $0.529$ at the primary cell and $0.549$ at $b=13$.
\item \textbf{Assumption dependence.} The verdicts of C1 to C4 are recomputed over the full grid $b\in\{1,2,5,10,13\}$, $C\in\{0.5,1,2\}$, $r\in\{1.5,2,3\}$, $\varepsilon\in\{0,0.02,0.05,0.10\}$, $L\in\{L_0/2,L_0,2L_0\}$. If the verdict of C1 is not constant over that grid, the certificate is recorded as assumption-driven on this record, and the sub-grid on which it passes is reported. This is a criterion about the chapter's own honesty rather than about the theorem, and it is expected to bind.
\end{enumerate}

\begin{table}[htbp]
\centering\small
\caption{Specification of the certificate table: one block of rows per $b\in\{1,2,5,10,13\}$, repeated over the mixing grid $(C,r)$ and the drift grid. It serves criteria C1 to C4 and reports the three deficit terms separately. No entry is filled here.}
\label{tab:usd-cert}
\begin{tabular}{@{}ll@{}}
\toprule
Column & Content \\
\midrule
$b$; $n$; $k$; $k/(n+1)$ & block layout and the realised nominal level \\
$\hat q$ & threshold, currency units per block, and as a multiple of the block score standard deviation \\
$1-\alpha-b_{n,k}(0.05)$ & finite-sample deficit, computed exactly \\
$(n+1)\bar\beta$; $\bar\beta/\delta'$ & coupling deficits, marginal and conditional, at the stated $(C,r)$ \\
$\hat\varepsilon$ & drift deficit: assumed value, or plug-in with its $L$ \\
Certified marginal & $1-\alpha-\hat\varepsilon-(n+1)\bar\beta$ \\
Certified conditional & $b_{n,k}(0.05)-\hat\varepsilon-\bar\beta/0.05$, with confidence $1-\delta-\delta'-n\bar\beta$ \\
C1 informative? & whether $\alpha+\hat\varepsilon+(n+1)\bar\beta<1$; with the grid count $93/180$ in the block header \\
C3 informative? & whether $\alpha-1/(n+1)-\hat\varepsilon-(n+1)\bar\beta>0$; grid count $2/180$ \\
$M$; breaches; $\hat\alpha$ & deployment blocks, breach count, realised frequency \\
Binomial $p$; bootstrap $p$ & one-sided tests of C1, with the detectable frequency at this $M$ \\
C1--C3 verdicts & pass / fail / not informative, one column per criterion \\
C4 record & consistent / conservative / below the certified conditional level; never pass or fail \\
\bottomrule
\end{tabular}
\end{table}

\figplan{Two panels on a shared date axis over the calibration and deployment windows. Upper, in currency units per block: the block scores plotted at the end date of each block, calibration and deployment blocks distinguished; the threshold $\hat q$ as a horizontal line from the first deployment block; breaches marked; gaps as shaded vertical strips, so that the half of the record spent on gaps is visible. Lower, in units of probability: the three deficit terms as a stacked area against $b$ on a secondary horizontal axis, at $C=1$, $r=2$ and the primary drift value, with the certified marginal level as the residual at the top and a rule at one half for criterion~C2.}{The certificate on the deployment window, and where its level goes. The upper panel tests criterion~C1 by eye and shows whether breaches cluster; the lower is the decomposition criterion~C5 examines, showing that the coupling term consumes the whole level at short block lengths, the finite-sample term at long ones, and that the drift term is the only one estimated rather than assumed. The deficit panel uses no data beyond $n$ and the assumed class.\label{fig:usd-cert}}

\section{Methods compared}\label{sec:usd-methods}

The certificate is one way of turning a record of hedging errors into a threshold. Three others are in use, and this section defines them in the form the comparison of Section~\ref{sec:usd-routes} needs. No theorem in this thesis ranks any of them against the certificate under Assumption~\ref{ass:mix}; the comparison is empirical, and the honest content of it is a table of stated levels beside realised breach frequencies.

\subsection{Block bootstraps}

The moving block bootstrap of \citet{kunsch1989} resamples the record $e_{1:N}$ in blocks of fixed length $\ell$: draw $\lceil N/\ell\rceil$ starting indices uniformly from $\{1,\dots,N-\ell+1\}$ with replacement, concatenate and truncate to length $N$ \citep[\S2.3, eq.~(2.7)]{kunsch1989}. The stationary bootstrap of \citet{politis1994} replaces the fixed length by independent geometric lengths of mean $1/p$ and wraps the record circularly, so the resampled series is stationary conditionally on the data. In either case the statistic of interest, here the $(1-\alpha)$-quantile of the block score, is recomputed on each pseudo-series.

For definiteness the stationary-bootstrap threshold used here is built on the block scores $S_1,\dots,S_{n'}$ of contiguous, unseparated blocks ($n'=\lfloor N/b\rfloor$, about twice the certificate's $n$): pseudo-samples of length $n'$ are assembled from segments with independent $\mathrm{Geometric}(p)$ lengths started at uniform indices modulo $n'$, the $(1-\alpha)$ empirical quantile $\hat q^{\,\flat}$ is taken on each, and the median of $R$ such values is the bootstrap threshold, their $[\alpha/2,1-\alpha/2]$ percentiles its uncertainty. The cube-root order in $p\asymp n'^{-1/3}$ is the bias--variance balance for the arithmetic mean under fixed-length blocking \citep[Cor.~3.1, p.~1226]{kunsch1989}, carried across to the stationary bootstrap's geometric length by analogy and not by any result read here, and used for want of a quantile-optimal rule with known constants. Two qualifications belong with it: K\"unsch's coefficient is the strong ($\alpha$-) one, not the $\beta$-coefficient of Assumption~\ref{ass:mix}, and his block count assumes the record divides exactly by the block length, so the $\lceil N/\ell\rceil$-with-truncation convention above is this chapter's and is not attributed to him.

Two constants the method leaves free are fixed here in advance. The replication count is $R=2000$, which puts the Monte Carlo standard error of a median quantile below a tenth of the spacing between adjacent order statistics of the pseudo-sample at every block length in the sweep, with the seed of Table~\ref{tab:appC-seeds}; and $1/p=\lceil n'^{1/3}\rceil$ evaluated at the $n'$ each block length produces, so that the rule is a function of the record and not a choice ($n'=25$ and $1/p=3$ at $b=10$ on a calibration year). The bootstrap is run on the calibration window only, so that its threshold and the certificate's see exactly the same data, which is what puts the comparison of Section~\ref{sec:usd-routes} on equal footing.

The theory of these methods is asymptotic and is developed in \citet{lahiri2003}: for smooth functionals of means of a $\beta$-mixing sequence the block bootstrap is consistent when $\ell\to\infty$ with $\ell/N\to0$, and at $\ell\asymp N^{1/3}$ it is second-order accurate. A quantile of the block score is not such a functional but an order statistic of a dependent sequence, whose Edgeworth expansion involves the density of the score law at the quantile and the long-run variance of $\ind{S\le q}$, both unknown, with correction coefficients depending on the mixing rate through constants that are not available. A statement of the form ``the bootstrap threshold covers the next block with probability at least $1-\alpha-\eta_N$'', with explicit $\eta_N$ and explicit constants, is therefore unavailable under Assumption~\ref{ass:mix} and is claimed neither here nor in Chapter~\ref{ch:certification}. No theorem ranks the two methods: the comparison below is empirical.

\gap{Citation locator needed: a chapter or theorem in \citet{lahiri2003} for the consistency of the block bootstrap under $\ell\to\infty$, $\ell/N\to0$ and for second-order accuracy at $\ell\asymp N^{1/3}$. The book has not been read and no locator is recorded for it; this is the same unlocated claim that Chapter~\ref{ch:certification} carries, and the two must be closed together. \citet{politis1994} carries no locator at all, the paper not having been obtained, so nothing beyond the descriptive definition above is attributed to it; the transfer of the fixed-block cube-root rule to the geometric-length rule is likewise uncited. \citet{kunsch1989} carries only the locators quoted above.}

The bootstrap uses the whole record rather than half of it, needs no gaps, and adapts to the dependence through $\ell$ or $p$ without an assumed class; its level, against that, is asymptotic and of unknown finite-sample accuracy. Which of the two is tighter on this record is an empirical question and is treated as one.

\begin{remark}[What the comparison is not]\label{rem:usd-notscp}
A negative finding in the literature concerns split block-\emph{permutation} conformal prediction, which \citet[p.~15]{stocker2025} report as failing to provide valid coverage even in simple stationary settings. That procedure is not the one certified here: it permutes blocks to construct a reference distribution, whereas the construction of Chapter~\ref{ch:certification} takes the order statistic of gap-separated block scores and pays for the dependence with an explicit coupling term. The negative finding does not transfer, and it is recorded here so that a reader does not import it. The nearest application of a conformal coverage bound in finance is \citet[Thm.~1]{schmitt2026}, whose coverage gap carries an unremovable factor two on a total-variation distance; the route of Chapter~\ref{ch:certification} carries the constant one on a Kolmogorov distance, and Chapter~\ref{ch:sharpness} shows that constant to be sharp. That paper reports regulatory breach-count statistics beside conformal coverage \citep[Table~1, p.~4]{schmitt2026}, which is the form Table~\ref{tab:usd-routes} follows.
\end{remark}

\subsection{Breach-count tests}

A threshold, however obtained, can be tested after the fact by counting breaches. Let $I_t:=\ind{S_t>q_t}$ be the breach indicator over $M$ deployment blocks with nominal breach probability $\alpha$, and $x:=\sum_tI_t$.

The proportion-of-failures test of \citet{kupiec1995} is the likelihood-ratio test of $H_0:\Prob(I_t=1)=\alpha$ against the unrestricted Bernoulli alternative,
\[
\mathrm{LR}_{\mathrm{uc}}:=-2\log\frac{(1-\alpha)^{M-x}\alpha^{x}}{(1-\hat\alpha)^{M-x}\hat\alpha^{x}},\qquad\hat\alpha:=x/M ,
\]
asymptotically $\chi^2_1$ under $H_0$ when the $I_t$ are independent. The test of \citet{christoffersen1998} adds a test of independence: with $n_{ij}$ the number of transitions from $I_{t-1}=i$ to $I_t=j$, and $\hat\pi_i:=n_{i1}/(n_{i0}+n_{i1})$, $\hat\pi:=(n_{01}+n_{11})/(n_{00}+n_{01}+n_{10}+n_{11})$,
\[
\mathrm{LR}_{\mathrm{ind}}:=-2\log\frac{(1-\hat\pi)^{n_{00}+n_{10}}\hat\pi^{n_{01}+n_{11}}}{(1-\hat\pi_0)^{n_{00}}\hat\pi_0^{n_{01}}(1-\hat\pi_1)^{n_{10}}\hat\pi_1^{n_{11}}}\ \sim\ \chi^2_1 ,
\]
and $\mathrm{LR}_{\mathrm{cc}}:=\mathrm{LR}_{\mathrm{uc}}+\mathrm{LR}_{\mathrm{ind}}$ is asymptotically $\chi^2_2$. Both are asymptotic in $M$, both presume at most first-order Markov dependence under the alternative, and both are diagnostics of a threshold already in use, not constructions of one. Section~\ref{sec:usd-routes} applies them to the certificate and bootstrap thresholds alike on the same deployment blocks. With $M$ of a few tens the asymptotic critical values are not to be trusted, so both are reported with a finite-sample $p$-value: exact binomial enumeration for the unconditional part, the block bootstrap of Section~\ref{sec:usd-ident} for the conditional part.

\gap{Citation locator needed: no page, equation or section locator may be quoted from \citet{kupiec1995} or \citet{christoffersen1998}, neither text having been obtained; the two statistics are stated above in the form in which they are used and are cited to the articles as wholes. The two properties of $\mathrm{LR}_{\mathrm{ind}}$ on which this chapter relies -- that it presumes a first-order Markov alternative and that its null distribution is asymptotic, so that it is a diagnostic and never a finite-sample guarantee -- are the intended use recorded for that article in this thesis's reading notes, and are stated above rather than attributed to a locator.}

\subsection{Regulatory backtests}

The Basel market-risk framework \citep{bcbs2019} contains two backtests relevant here. The first is the traffic-light count of value-at-risk exceptions: a bank that records too many exceptions in its one-year backtest is moved into a higher penalty zone and the multiplier applied to its market-risk capital charge rises accordingly, as \citet[\S1, p.~1]{schmitt2026} describes it, that being a secondary source and the only one read for this chapter. The length of the backtest window, the confidence level and the exception counts that separate the zones are boundary values of the framework and none of them is asserted here. The second is the profit-and-loss attribution test of the internal-models approach, comparing the risk-theoretical profit and loss of the desk's pricing model with the hypothetical profit and loss of the front-office system over the same days, by a Spearman rank correlation and a Kolmogorov--Smirnov distance between the two empirical distributions, with zones defined by thresholds on the two statistics. It is the regulatory analogue of the certificate's question, a systematic hedging error being exactly a gap between model and realised profit and loss, but it is a two-sample comparison at fixed sample size with asymptotic critical values and says nothing about the next block.

The chapter reports both statistics, and the zone only as read from the text at the time of the run. No conclusion here depends on where the boundaries fall: the comparison of interest is between the attribution statistics of the certified policy and those of the alternatives of Section~\ref{sec:usd-routes}, a within-record ranking invariant to the boundaries.

\gap{Citation locator needed: the paragraph numbers of \citet{bcbs2019} for the traffic-light zones and for the profit-and-loss attribution thresholds, together with the four numerical boundary values the framework fixes -- the length of the backtest window, the confidence level of the value-at-risk measure, the exception counts separating the zones, and the two attribution thresholds. No reading note exists for that document, so none of the four is stated anywhere in this chapter; the two statistics are reported and the zone is read from the January 2019 text and the consolidated framework at the time of the run.}

\section{The blocked and unblocked routes on equal footing}\label{sec:usd-routes}

\subsection{What equal footing requires}

Chapter~\ref{ch:sharpness} establishes that blocking costs a factor $\sqrt b$ in effective sample size and that the blocked calibration-conditional rate is not minimax in any coverage formulation; what blocking buys is the exact $\Beta(k,n+1-k)$ law and a failure probability whose only dependence on the mixing class is one term, not a better rate. Its Remark~\ref{rem:sharp-constants} adds that the practical ordering at a given record length turns on which inequality is used on the unblocked side, the crossover sitting near two thousand two hundred periods on the printed accounting and near seven hundred under the refined variance proxy. A record of a few hundred to a few thousand daily observations is exactly what this transition affords, so the comparison is not settled by theory at these lengths and must be run.

Three conditions make it fair, and all three are fixed here.

\paragraph{Same object.} Every route produces a threshold for the same block score at the same $b$ on the same deployment blocks, differing only in how the calibration record becomes a threshold: the certificate takes the $k$-th order statistic of $n=\lfloor N/(2b)\rfloor$ gap-separated block scores; the unblocked route the empirical $(1-\alpha)$-quantile of the $n'=\lfloor N/b\rfloor$ contiguous scores, whole record, no gaps; the stationary bootstrap the median of $R$ pseudo-sample quantiles from those same $n'$ scores.

\paragraph{Same cost accounting.} The blocked deficit $(n+1)\beta(b)$ already contains the coupling that decouples the deployment block from the threshold, and the unblocked route pays nothing for it, so comparing them as printed would credit the unblocked route with something it has not bought. Following Chapter~\ref{ch:sharpness}, the deployment coupling is removed from the blocked side so that every stated deficit is a calibration-side quantity, and the charge $\beta(b_g)/\delta'$ the unblocked route would owe for the same deployment statement is displayed in its own column rather than silently omitted.

\paragraph{Same sharpness of inequality, or an explicit note that it is not.} The blocked level rests on an exact Beta quantile, the unblocked level, in the form Chapter~\ref{ch:sharpness} proves, on Chebyshev. These are not of comparable sharpness, and a conclusion drawn from setting them side by side is an artefact of the inequalities. The table therefore carries three unblocked columns, each labelled with its status on the face of the table: the Chebyshev bound, which is proved; the sub-Gaussian size an exponential inequality of Fuk--Nagaev type would deliver, which Chapter~\ref{ch:sharpness} reports as an estimate; and the realised breach frequency, which is neither.

\subsection{The record lengths this transition affords}

In Theorem~\ref{thm:block-b} the record length $N$ is the calibration record on which a threshold is computed, not a span of data availability, so each of the three lengths is named with the record it denotes. The first is the calibration year, 2~January to 31~December 2022, about $250$ periods, the record of the certificate test of Section~\ref{sec:usd-cert}. The second is the whole predecessor record, 2~January 2019 to 30~June 2023, about eleven hundred periods, the longest on which the policy's predecessor fit exists on every date. The third runs from the successor series start of 2~April 2018 to 31~December 2024, about seventeen hundred periods; it is reachable only because the certified policy is frozen at the end of $\Dfit$ and so keeps producing block scores after the predecessor market ends.

The comparison runs over two axes, as criterion~R2 requires: the blocked route against each of the two unblocked accountings, at each of the three lengths. Evaluating $g_{\mathrm{cal}}(b)=1/(2\sqrt{n+2})+nCb^{-r}$ over \emph{feasible} block lengths only, against the Chebyshev bound Chapter~\ref{ch:sharpness} proves and the sub-Gaussian size it reports as an estimate, at $C=1$, $r=2$, $\delta=0.05$:

\begin{center}\small
\begin{tabular}{@{}rccc@{}}
\toprule
$N$ & blocked, minimised over feasible $b$ & unblocked Chebyshev (proved) & unblocked sub-Gaussian (estimate) \\
\midrule
$250$ & $0.2040$ & $0.7526$ & $0.3232$ \\
$1100$ & $0.1313$ & $0.3588$ & $0.1541$ \\
$1700$ & $0.1173$ & $0.2886$ & $0.1240$ \\
\bottomrule
\end{tabular}
\end{center}

Against the proved bound the blocked route is the smaller at all three lengths, by factors of $3.7$, $2.7$ and $2.5$. Against the sub-Gaussian estimate it is the smaller at all three on the printed accounting, but under the refined variance proxy of Chapter~\ref{ch:sharpness}, whose unblocked value is $3.879/\sqrt N$, the unblocked side reads $0.2454$, $0.1170$ and $0.0941$, so the blocked route is the smaller at $250$ periods and the larger at the other two. The ordering is therefore decided by the status of the inequality used on the unblocked side and not by the record length -- a sharper statement than a crossing, and the one this chapter reports.

The feasibility constraint of Section~\ref{sec:usd-cert} costs the blocked route something only at the shortest length: minimising $g_{\mathrm{cal}}$ over $b\le13$ gives $0.2040$ at $b=13$ against $0.1909$ at the unconstrained minimiser $b=17$, a penalty of $6.9$ per cent, while at $N=1100$ and $N=1700$ the feasible ceilings $61$ and $94$ lie above the minimisers $32$ and $38$ and the constraint does not bind. It moves neither crossing, which remain at about $2.2\times10^3$ and $1.0\times10^5$ periods. The left panel of Figure~\ref{fig:usd-routes} accordingly plots the constrained minimisation and marks the unconstrained crossings.

\subsection{Acceptance criteria, stated in advance}

\begin{enumerate}[label=\textbf{R\arabic*.},leftmargin=2.2\parindent]
\item \textbf{Reporting.} For each route: the threshold in currency units per block, the stated level with its status (proved finite-sample, proved asymptotic, estimated, or none), the realised breach frequency, and the two breach-count statistics with their finite-sample $p$-values. No route is declared to dominate. The criterion is satisfied by completeness, and is stated as a criterion so that a selective report is a visible failure.
\item \textbf{Ordering against the theory.} At each of the three record lengths the ordering of the blocked and unblocked deficits is compared with the ordering the crossover predicts under each of the two unblocked accountings. Agreement is recorded; disagreement is evidence about the constants in those inequalities and not about the theorem, neither the Chebyshev bound nor the sub-Gaussian estimate being claimed tight.
\item \textbf{Realised coverage is not a bound.} A route's realised breach frequency is reported without being set against that route's stated level as if a match confirmed it: a bootstrap threshold covering at the nominal level has not thereby acquired a finite-sample guarantee, and a certificate covering far above its stated level has not thereby been shown loose on the next record. The criterion is the absence of such inferences from the text.
\end{enumerate}

\begin{table}[htbp]
\centering\small
\caption{Specification of the route-comparison table: the four routes crossed with the three record lengths and $b\in\{5,10,13\}$, the deployment blocks common to every row. It serves criteria R1 and R2 and is the empirical content of the comparison Chapter~\ref{ch:sharpness} leaves open. No entry is filled here.}
\label{tab:usd-routes}
\begin{tabular}{@{}ll@{}}
\toprule
Column & Content \\
\midrule
Route & certificate (blocked, gaps); sample quantile (unblocked); stationary bootstrap; certificate at $b^*$ where feasible \\
$N$; $b$; $n$ or $n'$ & record length, block length, number of calibration scores \\
$\hat q$ & threshold, currency units per block \\
$\hat q$ ratio & threshold as a multiple of the certificate's threshold at the same $b$ \\
Stated level & the level the route states \\
Status & proved finite-sample / proved asymptotic / estimated / none \\
Calibration deficit & blocked: $1-\alpha-b_{n,k}(\delta)+n\bar\beta$; unblocked: Chebyshev, and sub-Gaussian estimate \\
Deployment charge & $\beta(b_g)/\delta'$, shown for the unblocked routes, zero for the certificate \\
$M$; breaches; $\hat\alpha$ & deployment blocks, breach count, realised frequency \\
$\mathrm{LR}_{\mathrm{uc}}$; exact $p$ & proportion-of-failures test with the exact binomial $p$-value \\
$\mathrm{LR}_{\mathrm{cc}}$; bootstrap $p$ & conditional-coverage test with a block-bootstrap $p$-value \\
Attribution & Spearman correlation and Kolmogorov--Smirnov distance of the attribution test, with the zone \\
\bottomrule
\end{tabular}
\end{table}

\figplan{Two panels. Left, both axes logarithmic, deficit in units of probability against record length $N$ from $10^2$ to $10^6$: the blocked calibration-side deficit minimised over admissible $b$, the unblocked Chebyshev bound and the unblocked sub-Gaussian estimate, at $C=1$, $r=2$, $\delta=0.05$, with the two crossings marked and vertical rules at the three record lengths this transition affords. Right, threshold in currency units per block against $b$ over the admissible range: the certificate threshold, the unblocked sample quantile and the median stationary-bootstrap threshold with its percentile band, all on the calibration year, with a rule at the largest realised deployment block score.}{The two routes at the record lengths this transition affords. The left panel carries no data and tests criterion~R2 by showing which side of each crossing the record falls on; the right is its empirical counterpart and tests whether the ordering of stated deficits is reproduced in the ordering of thresholds, which it need not be, a deficit being a statement about coverage and a threshold about currency.\label{fig:usd-routes}}

\datanote{The table of Section~\ref{sec:usd-routes} and the left panel of Figure~\ref{fig:usd-routes} use no market data: they evaluate $g_{\mathrm{cal}}(b)=1/(2\sqrt{n+2})+nCb^{-r}$ with $n=\lfloor N/(2b)\rfloor$, the Chebyshev form $\sqrt{(\tfrac12+4C\zeta(r))/(N\delta)}$, and the sub-Gaussian size $\sqrt{2\sigma_r^2\log(2/\delta)/N}$ with $\sigma_r^2=\tfrac14+2C\zeta(r)$, all at $C=1$, $r=2$, $\delta=0.05$, exactly as Chapter~\ref{ch:sharpness} specifies them; the third of these is an estimate and not a proved bound and the figure legend says so. The blocked column of the table minimises $g_{\mathrm{cal}}$ over integers $b$ with $\lfloor N/(2b)\rfloor\ge9$, which is the feasibility constraint of Section~\ref{sec:usd-cert}; the refined unblocked values replace $\zeta(2)=1.644934$ by the $S_2(1)=0.894934$ of Chapter~\ref{ch:sharpness}. The right panel requires the block scores of the certified policy over the calibration year, and therefore the whole of the data chain of Section~\ref{sec:usd-cert}.}

\section{What would falsify each result}\label{sec:usd-falsify}

A test is worth running only if some outcome would be taken as refuting the claim. The list below fixes those outcomes in advance, one entry per claim, and is deliberately separated from the criteria so that the criteria cannot be reinterpreted after the fact.

\begin{enumerate}[label=\textbf{F\arabic*.},leftmargin=2.2\parindent]
\item \textbf{The decomposition.} Theorem~\ref{thm:ident} says the transition map is a strike shift by the spread plus the forward basis, followed by the time change, followed by a perturbation in a two-dimensional family. It is falsified on this record by a containment rate that fails criterion~V1 with the failures \emph{gross} -- a median absolute normalised excess above one, so that the observation sits more than a full half-width outside the interval -- and \emph{dispersed} across expiries and quarters rather than concentrated in one episode. Gross dispersed failure cannot be attributed to a stale basis-volatility estimate or to a single stressed quarter, and would say that the contamination does not lie where the theorem puts it. It can, however, be produced by an estimate that is not stale but \emph{biased}: a volatility risk premium is a level effect, and a $\hat\eta$ too small by a constant factor fails every pair, everywhere, grossly. A third condition is therefore attached. The implied ratio $\eta^{\mathrm{imp}}/\hat\eta$ of Table~\ref{tab:usd-ident-excess} must \emph{not} be concentrated near a single constant; if it is, the finding is a level error in $\hat\eta$, is read under F2, and F1 does not fire.
\item \textbf{The exactness at the money.} Theorem~\ref{thm:ident-two}(i) is an identity in the basis scale and is not falsifiable by containment evidence. What such evidence can falsify is its \emph{input}, that the pricing-measure basis volatility is well proxied by the trailing realised one. A V1 failure whose excesses are narrow, with the basis volatility needed for containment systematically above $\hat\eta$, indicts the proxy and not the theorem; the reading is available only because criterion~V5 measures the same gap independently.
\item \textbf{The skew statement.} The first-order skew statement of Theorem~\ref{thm:ident} is falsified by a V2 failure with V1 passing, but only where the fitted basis scale is small: at plausible parameters the second-order term in the skew entry matches the first and can reverse its sign, so elsewhere a V2 failure measures the remainder. The realised scale $\hat\eta/a$ is reported by expiry and quarter, and the falsification reading holds only where it is below one fifth.
\item \textbf{The usefulness claim.} The claim that the interval is operationally useful is falsified by a failure of criterion~V3, a median relative half-width above one quarter on the majority of expiries. That outcome is compatible with every theorem of Chapter~\ref{ch:identifiability} being true, and is the outcome its own illustrative parameters would produce: it falsifies the thesis's practical claim and not its mathematical one, and the distinction is maintained in the reporting.
\item \textbf{The improvement claim.} The claim that the construction of Chapter~\ref{ch:identifiability} beats the market convention is falsified by a failure of criterion~V4 on the long-expiry arm, where the theory predicts the gain, with the block-bootstrap $p$-value agreeing with the sign test. A failure on the pooled arm alone is not decisive, because the short expiries carry the smallest contamination and the least room for improvement. If the grid truncates to $\{1,2\}$ years the long-expiry arm is empty; the improvement claim is then \emph{not tested}, F5 does not fire in either direction, and no pooled outcome may be substituted for it.
\item \textbf{The shared-basis assumption.} The assumption that one basis vector serves every expiry, which makes the deficit three rather than $3m$ and the book reserve additive in the form of Corollary~\ref{cor:ident-book}, is falsified by criterion~V6 recording an interquartile range above the stated threshold or per-expiry medians differing by more than it. That does not falsify Theorem~\ref{thm:ident}; it selects the per-expiry variant, changes the deficit count and changes the reserve of Chapter~\ref{ch:zar} from a net sum to a sum of absolute values.
\item \textbf{The certificate.} Theorem~\ref{thm:cov}(i) is falsified by a realised breach frequency significantly above the certified marginal deficit, at the primary cell and over a majority of the \emph{informative} cells of the assumption grid, with the excess not attributable to a single deployment block. One breach above the bound is not a refutation: the criterion is the binomial test of C1 with its bootstrap counterpart, and the requirement that the verdict survive the grid is criterion~C5.
\item \textbf{The certificate's assumptions.} Separately, the \emph{applicability} of the certificate to a record of this kind is falsified by criterion~C2: a certified level below one half at every feasible block length and every cell of the mixing grid says the certificate tells one nothing usable about a record of this length, whatever its realised coverage.
\item \textbf{The blocking comparison.} Nothing here is falsified by the unblocked route producing a tighter threshold or better realised coverage; Chapter~\ref{ch:sharpness} predicts exactly that beyond the crossover. Evidence against its accounting would be the blocked deficit exceeding the unblocked Chebyshev bound well below the crossing computed there, or falling below the sub-Gaussian estimate well above it, by a factor the stated constants cannot absorb.
\item \textbf{The design itself.} The experiment is invalidated, rather than any claim falsified, if the quoted-only rule leaves a grid too thin to support the stated denominators, or if $m^*$ lies outside the predecessor's quoted window so that $W_{\mathrm{ident}}$ is empty; then the strict arm is not run, the exchange arm is reported alone, and the chapter says so in place of a result. Two further outcomes are recorded rather than improvised later. If the expiry grid truncates, V4 and F5 are not evaluable and the remaining criteria are evaluated on the truncated grid at their stated thresholds. And if the predecessor volatility strip is refused altogether -- the one refusal Appendix~\ref{app:data} names as capable of stopping the first half of this chapter -- then Section~\ref{sec:usd-ident} stands unexecuted and, since Definition~\ref{def:usd-policy} fits that strip on every date, the certificate test does not survive either; what survives is the block-layout and deficit arithmetic, which uses no market data, and the chapter would report that alone and say so.
\end{enumerate}

Three outcomes are named here as \emph{not} falsifying anything, because each is easily mistaken for a refutation. A disagreement in sign between the option-implied and the historically estimated basis-rate correlation is evidence of a risk premium on basis volatility, which the theory of Chapter~\ref{ch:identifiability} neither asserts nor denies. A high containment rate on wide intervals is not confirmation, for the reason criterion~V3 is stated in Section~\ref{sec:ident-validation}. And a realised breach frequency far below the certified level is not evidence that the certificate is correct; it is evidence that it is conservative on this path, which is the quantity Chapter~\ref{ch:sharpness} treats.

\section{Threats to validity}\label{sec:usd-threats}

\subsection{The predecessor market was liquid; the South African one is not}

The dominant threat is external. This chapter tests the mathematics on a transition whose predecessor was one of the deepest option markets in the world. The South African predecessor market is not comparable in depth, its inter-dealer quotes are set against a collateral convention in a foreign currency, and the successor market did not exist at the time of the survey recorded in Chapter~\ref{ch:intro}. Three consequences follow, each carried into Chapter~\ref{ch:zar} rather than assumed away. The sampling error of the fitted predecessor triple is larger in a thin market, so the same theorem delivers a wider effective interval there at the same basis volatility, and no fit-residual diagnostic reported here transfers as a magnitude. The quoted-or-derived distinction, this chapter's defence against circularity, may not exist as a vendor flag there, in which case the corresponding test cannot be run at all. And the basis dynamics differ: the one published dynamic study of the South African spread models it with jumps at monetary-policy dates, which Chapter~\ref{ch:identifiability} records would break the cumulant computation at third order. A pass of every criterion here is evidence that the construction is sound, not that the numbers transport.

\subsection{Circularity in both directions}

A vendor successor surface before it is genuinely quoted is the predecessor surface shifted by the spread, and a vendor predecessor surface after trading has moved is the successor surface shifted back; testing a prediction against either would be circular, in opposite directions. The windows of Section~\ref{sec:usd-protocol} exclude both at the cost of a short strict arm, and the exchange arm exists because it rests on settlements from a market quoted by construction. The residual threat is that the vendor flag is itself unreliable, a point flagged as quoted resting on one stale contribution; nothing in the data resolves this, and the discarded count by window and expiry is the only visibility the chapter has.

\subsection{Composite surfaces, grids that move, and conventions that change}

The quality issues of composite volatility cubes, of strike grids that are re-centred as rates move, of mixed quoting conventions in the years when the predecessor forward was near zero, and of the contract-convention change that falls inside the window are recorded in Appendix~\ref{app:data} and are not restated here. Two of them shape the analysis and so belong in this list. The first is that every remedy leaves a residue, and the residue is largest exactly where the smile is most informative, in the wings, which is where criterion~V2 is measured. The second is that the remedies are all exclusions, so each of them thins the grid and moves the denominator of every rate reported in Section~\ref{sec:usd-ident}; that is why the discarded count is a column of Table~\ref{tab:usd-ident-summary} and not a footnote.

\subsection{The estimated input, and the measure it lives under}

The whole width of the identified interval rests on one estimated number, estimated under the physical measure while the theorem needs it under the pricing measure. The sampling error is small and is reported; the measure gap is not a sampling error, has no standard error and is not reduced by any window length. Criterion~V5 is the only instrument the chapter has for it and it measures a correlation rather than a volatility, so the gap in the \emph{level} remains unmeasured. This is the single largest threat to the quantitative content of Section~\ref{sec:usd-ident} and the reason the identified set is called an ambiguity set throughout, never an estimate with an error bar.

\subsection{Jumps at scheduled policy dates}

Successor forwards move in steps at scheduled policy dates. The chapter's rule -- expiries of one year and longer only, so that the step component is absorbed by the fit on both sides -- is an assumption and is gap-boxed in Section~\ref{sec:usd-protocol}. If the absorption is imperfect and asymmetric between the two benchmarks, the fitted triple carries a bias that the identification test would attribute to the basis. The diagnostic available is the fit residual by expiry: a residual that rises systematically at the shortest expiries in the grid is the signature of an unabsorbed step component, and it is reported for that purpose.

\subsection{Execution frictions, set to zero}

Both certified policies rebalance at every daily close, one in a benchmark in the middle of being replaced and the other, in addition, in an option market that is being born. Rebalancing in either carries bid-offer and slippage, and this construction charges neither: $e_t$ is computed from marks and not from fills. The consequence is signed. A study that charges no transaction cost understates the loss on every period by an amount scaling with the rebalancing frequency and the spread, so every block score, and therefore every threshold $\hat q$ here, is biased downward. The certificate is untouched, bounding the coverage of whatever score it is given; what is touched is the reading of $\hat q$ as a reserve, which is the reading Chapter~\ref{ch:zar} inherits. No spread data are requested in Appendix~\ref{app:data}, so the size of the bias is named and signed but not estimated.

\subsection{Multiplicity, and the single path}

Six identification criteria, five certificate criteria and three route criteria are evaluated, several per expiry. No multiplicity correction is applied and none is claimed: the criteria are reported jointly with the count passed and no subset is presented as the result. The deeper limitation is that all are computed on one path of one market; the block bootstrap measures sampling variability along that path and says nothing about paths that were not realised, as no procedure can.

\subsection{The policy is a convention, not a book}

The certified policy is the market convention applied to a declared synthetic portfolio, not any institution's policy on its real book. The certificate is indifferent, requiring only that the policy be fixed; the economic reading of the threshold is not, since the currency units of $\hat q$ are those of the declared portfolio and scale with its declared vega. Every threshold is therefore reported in currency units and as a multiple of the block-score standard deviation, only the second being comparable across books.

\section{Discussion}\label{sec:usd-discussion}

\subsection{What the chapter will have shown}

The chapter will have shown three things: whether a smile predicted from a predecessor market, a pair of curves and a published constant contains the successor smile the market later quoted, at a rate fixed in advance; whether the interval so produced is narrow enough to be used, which is an independent question and the one a desk asks; and whether a finite-sample coverage certificate deployed at the market convention delivers what it states. The first is a statement about a map, the prediction being an interval from an identification argument and not a forecast. The second may fail while the first passes, in which case the identification result bounds the error of the market convention without replacing it -- a smaller claim than the thesis would like, and a defensible one. The third carries the part most likely to generalise, which is not the verdict but the decomposition: that the coupling term consumes the whole level at short block lengths while the finite-sample term does so at long ones, and that the drift term is the only one estimated rather than assumed, is a statement about the certificate's structure that Chapter~\ref{ch:zar} inherits directly.

\subsection{What a failure would mean}

The chapter is designed so that each way of failing points somewhere. A gross, dispersed containment failure indicts the model class. A narrow failure with an implied basis volatility above the estimate indicts the estimate and, through it, the measure gap. A failure of the width criterion with the containment criterion passing indicts the usefulness of the result and not its truth. A certificate breach frequency above the bound at every cell of the assumption grid indicts the mixing hypothesis, which is the only hypothesis in Theorem~\ref{thm:cov} that cannot be checked from the data, since Chapter~\ref{ch:certification} shows the coefficient is not estimable distribution-free. And a certified level below one half at every feasible block length says that the record is too short for the certificate, which is a statement about records rather than about theorems and is the one most likely to be true.

The chapter's own arithmetic supports that last reading before any datum arrives. Over the pre-registered grid the certified marginal bound exceeds the parameter it bounds in $93$ of $180$ cells, and its two-sided companion can fail in $2$. The binding limitation at $N=250$ is therefore not Theorem~\ref{thm:cov} but the length of the record, and the criteria are written so that this shows as \emph{not informative} rather than banking as \emph{pass}.

\subsection{The comparison that is not a theorem}

Nothing here establishes that the certificate is better than the stationary bootstrap, or blocking better than not blocking: the two guarantees are of different kinds and their thresholds are comparable only as numbers on one record. Chapter~\ref{ch:sharpness} locates the length at which the unblocked route overtakes the blocked one; this chapter measures where the available record sits relative to it. That is the whole of the claim, and the criteria of Section~\ref{sec:usd-routes} exist to resist writing a ranking on the strength of one record.

\subsection{Connection to the converted book}

Chapter~\ref{ch:zar} takes three things from here and nothing else: the protocol, meaning a policy fixed in advance, windows justified from a cessation calendar and criteria stated before the record is seen; the deficit decomposition, with the warning that the South African record is shorter by an order of magnitude and monthly rather than daily, so that the feasibility constraint which merely displaced the preferred block length here will bind much harder there; and the finding, whichever way it goes, on whether the physical-measure basis volatility proxies its pricing-measure counterpart, the converted book's reserve being linear in that number with no other source. It takes no parameter estimate, for the reason set out in Section~\ref{sec:usd-threats}.

\datanote{No number in this chapter is computed from market data. Every table is a specification of columns and rows, every figure is a specification of axes and series, and every acceptance criterion is a threshold fixed before the corresponding data exist in the working set. The numbers that are computed here -- the feasibility ceiling $b\le13$; $b^*\approx17.8$; the sweep $n=125,62,25,12,9$ and its five values of $k/(n+1)$; $(n+1)Cb^{-r}=126$ at $b=1$; the grid counts $93/180$ and $2/180$; the Beta quantiles $0.7791$ and $0.7169$; the deficit table of Section~\ref{sec:usd-routes}; the binomial detection thresholds at $M=12$ and $M=25$; and $1/p=3$ at $n'=25$ -- use only the record length, the level $\alpha$, the assumed mixing class and the worked parameters of Chapter~\ref{ch:identifiability}, and each is stated with its inputs. The datasets named are those of Appendix~\ref{app:data}, whose unresolved identifiers are resolved in the single terminal session recorded there; no criterion, window or threshold here depends on an unresolved identifier.}
